%=====================================================================
\chapter{Command Reference}\label{app:commands}
%=====================================================================
\normalfigures

Organised by task rather than by tool, because that is how the question
arrives. Keep this open beside the labs.

\section{Is This Machine Capable?}

\begin{lstlisting}[language=bash]
grep -c -E '\<(vmx|svm)\>' /proc/cpuinfo   # VT-x / AMD-V present
grep -o -m1 -E '\<(ept|npt)\>' /proc/cpuinfo   # two-dimensional paging
ls /sys/kernel/iommu_groups | wc -l        # IOMMU on (Chapter 4)
cat /sys/module/kvm_intel/parameters/nested    # nesting (or kvm_amd)
sudo virt-host-validate qemu               # everything, with verdicts
systemd-detect-virt                        # am I already a guest?
sudo dmidecode -s system-manufacturer      # which hypervisor, if so
lscpu | grep -E 'Socket|Core|NUMA|Thread'  # topology (Chapter 6)
\end{lstlisting}

\section{Virtual Machines: \texttt{virsh}}

\begin{lstlisting}[language=bash]
# --- lifecycle --------------------------------------------------------
virsh list --all                     # every guest and its state
virsh start NAME                     # power on
virsh shutdown NAME                  # ask the guest to stop (ACPI)
virsh destroy NAME                   # pull the power. Not the same thing
virsh undefine NAME --remove-all-storage --snapshots-metadata
virsh console NAME                   # serial console; Ctrl-] to leave
virsh autostart NAME                 # start with the host

# --- inspecting -------------------------------------------------------
virsh dominfo NAME                   # the allocation
virsh dumpxml NAME                   # the whole definition
virsh domstats NAME --vcpu --balloon --block --interface
virsh dommemstat NAME                # balloon, swap_in, swap_out (Ch 6)
virsh domdirtyrate-calc NAME --seconds 10   # then domstats --dirtyrate
virsh domblklist NAME                # which files back which disks

# --- changing ---------------------------------------------------------
virsh setvcpus NAME 4 --config --maximum
virsh setvcpus NAME 4 --config
virsh setmaxmem NAME 8G --config
virsh setmem NAME 8G --config
virt-xml NAME --edit --network model=virtio      # Chapter 7
virt-xml NAME --edit target=vda --disk bus=virtio

# --- snapshots (Chapter 12) -------------------------------------------
virsh snapshot-create-as NAME tag "why I took it"
virsh snapshot-list NAME
virsh snapshot-revert NAME tag --running
virsh snapshot-delete NAME tag

# --- migration (Chapter 17) -------------------------------------------
virsh migrate --live --verbose --persistent --undefinesource \
      NAME qemu+ssh://DEST/system
virsh migrate --live --auto-converge NAME qemu+ssh://DEST/system
virsh migrate-setspeed NAME 1000     # MiB/s cap
virsh domjobinfo NAME                # progress of a running migration
\end{lstlisting}

\section{Virtual Disks: \texttt{qemu-img}}

\begin{lstlisting}[language=bash]
qemu-img create -f qcow2 disk.qcow2 20G
qemu-img create -f qcow2 -b base.qcow2 -F qcow2 child.qcow2   # Ch 12
qemu-img info disk.qcow2
qemu-img info --backing-chain disk.qcow2     # the whole chain
qemu-img convert -O raw disk.qcow2 disk.raw  # also vmdk, vhdx, vpc
qemu-img resize disk.qcow2 +10G
qemu-img check disk.qcow2
qemu-img snapshot -l disk.qcow2              # internal snapshots
du -h disk.qcow2 ; ls -lh disk.qcow2         # actual vs apparent size
\end{lstlisting}

\section{Containers: \texttt{docker}}

\begin{lstlisting}[language=bash]
# --- images (Chapter 9) -----------------------------------------------
docker build -t name:tag .
docker history name:tag --format '{{.Size}}\t{{.CreatedBy}}'
docker image inspect name:tag --format '{{json .RootFS.Layers}}'
docker save name:tag | tar -t | head          # the layers, as files
docker pull repo@sha256:...                   # pin a digest, not a tag

# --- running ----------------------------------------------------------
docker run --rm -it image sh
docker run --init image            # a real PID 1 (Chapter 8)
docker stats --no-stream
docker inspect -f '{{.HostConfig.Privileged}}' CONTAINER

# --- hardening (Chapter 19) -------------------------------------------
docker run --user 10001:10001 --cap-drop=ALL \
  --security-opt=no-new-privileges --read-only --tmpfs /tmp \
  --pids-limit 100 --memory 256m --cpus 0.5 image

# --- the kernel primitives underneath (Chapter 8) ---------------------
ls -l /proc/self/ns/                          # this process's namespaces
unshare --pid --mount --uts --ipc --net --fork --mount-proc CMD
systemd-cgls                                  # the cgroup tree
cat /sys/fs/cgroup/CGROUP/cpu.stat            # throttled_usec
cat /sys/fs/cgroup/CGROUP/memory.max
\end{lstlisting}

\section{Kubernetes: \texttt{kubectl}}

\begin{lstlisting}[language=bash]
kubectl get pods -o wide --all-namespaces
kubectl describe pod NAME | sed -n '/Events/,$p'   # why it will not run
kubectl get events --sort-by=.lastTimestamp | tail -20
kubectl top nodes ; kubectl top pods               # needs metrics-server
kubectl get pod NAME -o custom-columns=\
NAME:.metadata.name,QOS:.status.qosClass           # Chapter 10
kubectl describe node NAME | sed -n '/Allocated/,$p'   # requests
kubectl set resources deployment NAME --requests=cpu=500m,memory=2Gi
kubectl rollout status deployment NAME
kubectl rollout undo deployment NAME
kubectl logs -f POD --previous                     # the run that crashed
\end{lstlisting}

\section{Networking}

\begin{lstlisting}[language=bash]
# --- namespaces and virtual interfaces (Chapter 13) -------------------
ip netns add NAME ; ip netns exec NAME CMD
ip link add veth0 type veth peer name veth1
ip link add br0 type bridge ; ip link set veth0 master br0
bridge fdb show br br0                  # the switch's MAC table
bridge link show

# --- overlays ---------------------------------------------------------
ip link add vx0 type vxlan id 100 remote IP local IP dstport 4789 dev eth0
ip -d link show vx0                     # confirm the VNI
ip link set vx0 mtu 1450                # the arithmetic of Chapter 13

# --- diagnosis --------------------------------------------------------
ping -M do -s 1472 HOST                 # find the real path MTU
tracepath HOST                          # reports the MTU at each hop
tcpdump -ni IFACE udp port 4789         # watch the encapsulation
ethtool -k IFACE | grep -i offload      # is VXLAN offloaded?
ss -tuln                                # what is listening
tc qdisc add dev IFACE root netem delay 150ms    # simulate latency
tc qdisc del dev IFACE root
\end{lstlisting}

\section{Storage}

\begin{lstlisting}[language=bash]
# --- thin pools (Chapter 12) ------------------------------------------
pvcreate /dev/sdX ; vgcreate vg /dev/sdX
lvcreate -L 100G --thinpool pool vg
lvcreate -V 50G --thin -n vol vg/pool
lvs vg -o lv_name,lv_size,data_percent,metadata_percent
fstrim -v /mountpoint                   # return freed blocks to the pool

# --- measurement ------------------------------------------------------
iostat -x 2                             # aqu-sz and await are the ones
fio --name=t --rw=randread --bs=4k --size=1G --direct=1 \
    --runtime=30 --time_based --output-format=json
dd if=/dev/zero of=/tmp/t bs=1M count=2048 oflag=direct
lsblk -t                                # queue depths and scheduler
\end{lstlisting}

\section{Measurement}

\begin{lstlisting}[language=bash]
# --- from inside a guest (Chapter 18) ---------------------------------
mpstat 2 5                              # %steal is ready time, from inside
vmstat 2 5                              # si/so are swap; r is run queue
top                                     # %st in the CPU line

# --- from the host ----------------------------------------------------
perf stat -e 'kvm:kvm_exit' -p PID -- sleep 10       # VM exits (Ch 4)
perf record -e 'kvm:kvm_exit' -p PID -- sleep 10 ; perf script
virsh domstats --vcpu                   # wait and delay per vCPU
sar -u -r -b 2 5                        # cpu, memory, block, together

# --- generating load --------------------------------------------------
stress-ng --cpu 4 --timeout 60s --metrics-brief
stress-ng --vm 2 --vm-bytes 80% --vm-keep --timeout 60s
iperf3 -s -D ; iperf3 -c HOST -t 20 -P 4
\end{lstlisting}

\section{Shell Idioms Used in the Labs}

If any of these are unfamiliar, this is the section the prerequisite self-check
was pointing at.

\rowcolors{2}{white}{lightbg}
\begin{longtable}{@{}L{5.4cm}L{9.6cm}@{}}
\toprule
\rowcolor{primary!15}
\textbf{Idiom} & \textbf{What it does} \\
\midrule
\endhead
\texttt{set -euo pipefail} & Stop on the first error, on an unset variable, and
on a failure anywhere in a pipeline. Every lab script begins with it \\
\texttt{VAR=\$(command)} & Capture a command's output into a variable \\
\texttt{cmd \textbar{} awk '/pat/\{print \$2\}'} & Print the second field of
lines matching a pattern --- how the labs extract an address \\
\texttt{cmd \textbar{}\textbar{} true} & Do not let a failure here stop the
script; used around teardown \\
\texttt{name () \{ \ldots{} \}} & Define a function. The labs use these to keep
lines under 76 columns \\
\texttt{cmd <<'EOF' \ldots{} EOF} & A here-document. Quoted \texttt{'EOF'} means
nothing inside is expanded --- important for YAML and for anything containing
\texttt{\$} \\
\texttt{\&} then \texttt{wait} & Run in the background, then wait for all
background jobs. Used to make several guests busy at once \\
\texttt{for x in \$(seq 1 4); do \ldots{} done} & Repeat with a counter \\
\bottomrule
\end{longtable}
